\documentclass[UTF8,a4paper,11pt,fontset=fandol]{ctexart}
\usepackage{geometry}
\geometry{margin=2.2cm}
\usepackage{amsmath,amssymb,longtable,booktabs,array,hyperref,xcolor,enumitem}
\hypersetup{colorlinks=true,linkcolor=blue,urlcolor=blue}
\setlist[itemize]{leftmargin=1.3em,itemsep=0.2em}
\definecolor{notegray}{RGB}{246,248,251}
\title{P01. A Priority-Aware Replanning and Resequencing Framework for Coordination of Connected and Automated Vehicles\\中英双语博士精读笔记}
\author{Behdad Chalaki; Andreas A. Malikopoulos}
\date{2021 · arXiv preprint}
\begin{document}
\maketitle

\noindent\textbf{Theme / 主题：} Scheduling/resequencing + optimal control\\
\textbf{Method family / 方法族：} Priority-aware resequencing + decentralized optimal control / 优先级重排序与分散最优控制\\
\textbf{PDF pages / 页数：} 6 \quad
\textbf{Relevance / 相关度：} 94/100\\
\textbf{Source / 来源：} \url{https://arxiv.org/abs/2109.05573}

\section{论文全景速览与核心贡献 / High-Level Overview \& Core Contributions}
\subsection{研究背景与痛点 / Background \& Pain Points}
这篇论文位于“Scheduling/resequencing + optimal control”方向。对你的 JSSP-based Intersection Management 课题，它回答的是高层调度、低层轨迹平滑、安全过滤或真实验证中的一个关键子问题：如何让车辆在路口少停、少冲突、少能耗，同时保持在线可执行。

The paper belongs to the theme of Scheduling/resequencing + optimal control. For a JSSP-based intersection-management thesis, it illuminates one of the key layers: scheduling, smoothing, safety filtering, or real-world validation.

\textbf{Pain point / 痛点：} Assumes connected automated vehicles, known paths, and simplified dynamics; the paper is closer to resequencing theory than to learned graph scheduling.

\subsection{核心科学问题 / Core Scientific Question}
当车辆到达扰动使原始通行序列不再高效或可行时，如何在不牺牲安全约束的前提下快速重排序并生成可执行轨迹？

How can an intersection manager resequence disturbed arrivals while preserving safety and low-level trajectory feasibility?

\subsection{科学贡献 / Scientific Contributions}
\begin{itemize}
\item 方法贡献 (method contribution)：Event-driven replanning and priority-aware resequencing layered on decentralized CAV optimal control at a signal-free intersection.
\item 安全/避碰贡献 (safety contribution)：Uses rear-end and lateral safety constraints around conflict points, with replanning when disturbances change feasibility.
\item 平滑/停止避免贡献 (smoothing contribution)：The lower-level optimal-control trajectory is constructed so vehicles can conserve momentum rather than stop at the intersection when a feasible exit time exists.
\item 创新力度评判 (reviewer judgement)：强相关但中等创新：把 resequencing 明确嵌入 CAV 最优控制闭环，贡献主要在系统架构和事件触发逻辑。
\end{itemize}


\subsection{核心概念对照 / Core Concepts}
\begin{longtable}{p{0.24\linewidth}p{0.28\linewidth}p{0.40\linewidth}}
\toprule
中文概念 & English Concept & Role / 学术作用\\
\midrule
自动交叉口管理 & Autonomous Intersection Management & 把信号控制替换为车辆级调度与控制。 \\
轨迹平滑 & Trajectory Smoothing & 减少急加减速、停车和能耗峰值。 \\
安全约束 & Safety Constraints & 把碰撞风险写成距离、时间间隔或屏障函数。 \\
Priority-aware resequencing + decentralized optimal control & 优先级重排序与分散最优控制 & 该论文的核心方法族。 \\
\bottomrule
\end{longtable}

\section{理论方法论深度解构 / Methodology \& Theoretical Framework}
\subsection{问题数学建模 / Mathematical Formulation}
\textbf{双积分车辆模型与能量型最优控制 / Double-integrator vehicle model with energy-like optimal control}
\[
\min_{u_i(\cdot),\,t_i^f}\; J_i=\frac{1}{2}\int_{t_i^0}^{t_i^f}u_i(t)^2\,dt,\quad \dot p_i(t)=v_i(t),\quad \dot v_i(t)=u_i(t)
\]
\begin{longtable}{p{0.16\linewidth}p{0.25\linewidth}p{0.25\linewidth}p{0.25\linewidth}}
\toprule
Symbol / 符号 & 含义 & Meaning & Role / 建模作用\\
\midrule
i & 车辆索引 & vehicle index & 把多车协调拆成单车最优控制子问题 \\
p\_i(t) & 沿路径位置 & path-coordinate position & 描述车辆是否接近/离开冲突区 \\
v\_i(t) & 速度 & speed & 连接安全间距、通行时间与停止避免 \\
u\_i(t) & 加速度控制 & acceleration control & 优化变量，平滑性和能耗代理 \\
t\_i\textasciicircum{}f & 车辆离开控制区时间 & terminal exit time & 高层排序传给低层控制器的关键接口 \\
\bottomrule
\end{longtable}

\subsection{核心算法架构 / Core Algorithm Layout}
高层监测车辆进入、延误和可行性破坏；当扰动出现时，按照优先级规则更新局部通行序列，再把目标离开时间传给每辆车的解析/数值最优控制器。低层控制器检查速度、加速度、追尾和横向冲突约束是否可执行。

The upper layer detects event-triggered infeasibility, resequences vehicles under priority rules, and sends target exit times to decentralized optimal controllers that enforce motion and safety constraints.

\subsection{收敛性、稳定性或实时性 / Convergence, Stability, Real-Time Feasibility}
实时性来自局部事件触发与低维车辆模型；安全性依赖硬约束而非学习策略。它不证明全局最优，但把全局重排问题约束在小窗口内，适合在线运行。

Real-time behavior comes from event-triggered local resequencing and low-dimensional dynamics; safety is constraint-based, while global optimality is intentionally relaxed.

\subsection{潜在假设与边界条件 / Assumptions \& Boundary Conditions}
\begin{itemize}
\item 所有车辆为联网自动驾驶车辆 (connected automated vehicles, CAVs)，且路径已知。
\item 车辆动力学可由简化双积分模型 (double-integrator model) 充分描述。
\item 通信、定位和控制执行误差较小，扰动可被及时探测。
\end{itemize}


\section{实验验证与结果评判 / Experimental Validation \& Result Evaluation}
\textbf{Baselines \& Environment / 基准与环境：} 通常与固定先来先服务/不重排序 (FCFS / no resequencing) 策略比较，重点观察扰动后恢复效率。

\textbf{Validation / 验证：} Numerical simulation of signal-free intersection coordination under resequencing and replanning.

\textbf{Metrics / 指标：} 通行时间 (travel time)、延误 (delay)、停止次数 (number of stops)、控制能量代理 (control effort) 和安全约束违反。

\textbf{Ablation / 消融：} 论文更像机制验证而非完整消融；最应补充的是关闭优先级、关闭事件触发和扩大重排窗口的对比。

\textbf{Reviewer reading / 审稿式阅读提醒：} 阅读实验时不要只看平均值；应检查交通密度、车流方向、随机种子、失败案例和计算耗时。For doctoral reading, inspect density regimes, demand patterns, random seeds, failure cases, and computation time rather than only mean improvements.

\subsection{图表阅读线索 / Figure and Table Reading Cues}
PDF extraction detected approximately 11 figure references and 0 table references. Exact captions remain in the original PDF; the bilingual cues below tell you what to inspect first.

\begin{longtable}{p{0.24\linewidth}p{0.30\linewidth}p{0.36\linewidth}}
\toprule
图表名称 & Figure/Table Name & Reading focus / 阅读重点\\
\midrule
系统框架图 & system framework diagram & 先确认输入、决策层、控制层和安全约束之间的数据流。 \\
冲突关系图 / 通行顺序图 & conflict-relation or passing-order graph & 把节点、边类型和先后约束翻译成 JSSP 的作业、机器和析取约束。 \\
实验指标对比图表 & experimental metric comparison plots/tables & 重点看平均值之外的密度变化、失败场景、计算时间和方差。 \\
\bottomrule
\end{longtable}

\section{审稿人视角：批判性思考与未来方向 / Critical Thinking \& Future Directions}
\subsection{论文局限性 / Limitations \& Weaknesses}
\begin{itemize}
\item 对全 CAV 环境依赖强，混合交通中人类驾驶车辆 (HDVs) 会破坏可行性假设。
\item 优先级规则可能引入公平性问题，长期低优先级车流需要饥饿避免机制。
\item 简化动力学忽略执行器延迟、轮胎约束和感知误差累积。
\end{itemize}


\subsection{博士课题拓展点 / Research Extensions for PhD}
\begin{itemize}
\item 把固定优先级扩展为学习型优先级函数 (learned priority function)，并加入公平性正则。
\item 将重排序窗口映射为 JSSP machine queue，用图神经网络 (GNN) 预测候选交换收益。
\item 在低层加入鲁棒 MPC 或 CBF 安全过滤器，处理通信延迟和轨迹预测误差。
\end{itemize}


\subsection{如何服务当前课题 / Use in This Project}
Directly supports this project's hierarchy: a high-level sequence or schedule provides target times, and a low-level controller checks executable motion.

\section{交互式可视化与 PDF 排版蓝图 / Interactive Visualization \& PDF Layout Blueprint}
\textbf{动态参数 / Parameter：} 重排序窗口大小 / Resequencing window size, range [2, 12] veh.

窗口增大通常降低延误但提高计算量，并可能增加局部公平性风险。

A larger window usually reduces delay but increases computation and fairness risk.

\subsection{HTML 结构化布局建议 / HTML Structuring}
\begin{itemize}
\item 顶部固定 metadata bar：题名、作者、年份、主题、PDF 页数。
\item 左侧目录/section navigator，右侧为五大精读模块。
\item 公式卡片使用 <pre> 或 LaTeX block，紧跟符号表。
\item 审稿人批注用 reviewer-note block，博士拓展用 research-idea list。
\item 交互参数用 input[type=range] + canvas 曲线，打印 PDF 时隐藏交互控件但保留解释。
\end{itemize}


\section{来源锚点与逐节阅读路径 / Source Anchors \& Section-by-Section Reading Map}
\begin{longtable}{p{0.14\linewidth}p{0.76\linewidth}}
\toprule
Page & Extracted heading\\
\midrule
p.1 & I. INTRODUCTION \\
p.1 & II. M ODELING FRAMEWORK \\
p.1 & III. Finally, we provide simulation results in Section IV, and \\
\bottomrule
\end{longtable}

\paragraph{p.1 I. INTRODUCTION}定位问题背景、研究缺口和为什么现有保守/非协同方法不足。 / Frames the motivation, research gap, and limits of conservative or non-cooperative methods.

\paragraph{p.1 III. Finally, we provide simulation results in Section IV, and}检验方法是否真的改善效率、安全或能耗；要关注基准是否公平。 / Tests efficiency, safety, or energy claims; inspect whether baselines are fair.

\paragraph{p.1 II. M ODELING FRAMEWORK}给出算法流程和求解机制，应重点追踪输入、输出和安全接口。 / Gives the algorithmic mechanism; track inputs, outputs, and safety interfaces.

\paragraph{p.3 III. A PRIORITY-AWARE RESEQUENCING}作为局部论证段落阅读，关注它如何服务主问题。 / Read as a local argument and ask how it supports the main question.

\paragraph{p.5 IV. SIMULATION RESULTS}检验方法是否真的改善效率、安全或能耗；要关注基准是否公平。 / Tests efficiency, safety, or energy claims; inspect whether baselines are fair.

\paragraph{p.6 V. CONCLUDING REMARKS AND DISCUSSION}总结贡献和未来方向，也常暴露作者承认的边界条件。 / Summarizes contributions and often reveals author-recognized boundaries.


\end{document}
